%!TEX root = ../template.tex
%%%%%%%%%%%%%%%%%%%%%%%%%%%%%%%%%%%%%%%%%%%%%%%%%%%%%%%%%%%%%%%%%%%%
%% chapter1.tex
%% NOVA thesis document file
%%%%%%%%%%%%%%%%%%%%%%%%%%%%%%%%%%%%%%%%%%%%%%%%%%%%%%%%%%%%%%%%%%%%

\typeout{NT FILE chapter1.tex}%

\chapter{Introduction}
\label{cha:introduction}

\section{Motivation}
\label{sec:motivation}

\ntindex[Motivation]{}

Modern computing is inherently dependent on large-scale distributed systems, where high availability and low latency are essential requirements. To meet these demands in the presence of network failures and temporary disconnections, the CAP theorem~\cite{cap-theorem} implies that strong consistency must often be relaxed in favor of availability. In this context, Conflict-Free Replicated Data Types (CRDTs) have emerged as a fundamental solution, providing Strong Eventual Consistency by construction~\cite{shapiro2011comprehensive}. CRDTs allow distributed replicas to be updated independently and concurrently, while guaranteeing that all replicas eventually converge to the same state without requiring synchronous coordination.

Despite these guarantees, implementing CRDTs correctly is challenging and prone to errors that can compromise data integrity. Consequently, formal verification techniques play an important role in increasing confidence in their correctness. Existing approaches can be divided into fully-automated verification and deductive verification. Automated tools based on Satisfiability Modulo Theories (SMT) solvers, such as VeriFx~\cite{porre2023masses}, support validation of convergence properties and the automatic discovery of counterexamples. Deductive verification platforms, such as Why3~\cite{why3}, offer stronger guarantees by enabling proofs of functional correctness through contracts and invariants.

Experience shows that neither approach is sufficient when used in isolation. Automated verification, while efficient, struggles to handle complex or recursive structures. Deductive verification, although more expressive, often requires manual effort. This highlights the need for a methodology that combines the strengths of both approaches to support the verification of modern CRDTs.

\section{Problem Definition}
\label{sec:problem}

\ntindex[Problem Definition]{}

Despite advances in automated verification, tools such as VeriFx face clear limitations when applied to CRDTs with more complex behavior. The main problem addressed in this dissertation is that purely automated techniques are often unable to guarantee correctness in scenarios that require inductive reasoning or the preservation of global state invariants~\cite{dina2022thesis}.

Previous work has shown that SMT-based solvers exhibit well-known limitations when reasoning about recursive definitions~\cite{why3}. First, SMT-based solvers have difficulty reasoning about recursive definitions, properties that depend on iterating over data structures of unbounded size are hard to prove without explicit induction. This issue arises, for example, in structures with recursive insertion logic, such as the Replicated Growable Array (RGA)~\cite{roh2011rga}, where determining where a new element belongs requires walking an unbounded number of existing elements, often leading to timeouts or inconclusive results.

Another limitation of automated verification is that it does not guarantee that the implemented conflict-resolution policy aligns with the intended semantics. Convergence alone cannot distinguish between different behaviors, such as add-wins or remove-wins. In addition, verification results may be unstable, as small, non-functional code changes can lead to unexpected verification failures. These limitations show that automation alone is insufficient for verifying modern CRDTs, motivating the use of deductive techniques that can reason about induction, enforce invariants, and provide stronger semantic guarantees.

\section{Goals and Contributions}
\label{sec:contributions}

\ntindex[Goals and Contributions]{}

The contributions of this work are threefold. First, it introduces a dedicated language for development of CRDTs,
together with a source-to-source compiler that automatically derives both a VeriFx and a WhyML implementation from a
single specification, removing the need to author and maintain two independent formalizations by hand. Second, it delivers
a library of ten CRDTs specified in this language, nine of which are verified in both backends, covering state-based and
operation-based counters, sets, a generic map, and the Replicated Growable Array, together with the translation patterns
required to generate well-integrated code for both tools. Third, an experimental evaluation compares the specification
effort required from the programmer against the code generated for each target. For the Replicated Growable Array, the
VeriFx proof times out, while the Why3 proof, guided by inductive lemmas, establishes that every pair of operations except
two insertions commutes. Together, these results illustrate the practical value of combining both verification paradigms.
The results of this work were partially presented, as a paper, in INForum 2026, the 17th Portuguese Symposium in
Computer Science. The source code and all the examples used are available online~\footnote{Online at: https://github.com/JoaoGarcao/master-thesis}.

Within this context, the specific goals of the dissertation are as follows:

\begin{itemize}
    \item To design a single specification language capable of expressing both state-based (CvRDT) and operation-based
    (CmRDT) CRDTs.
    \item To formalize a source-to-source compilation pipeline that translates this language into both VeriFx and WhyML,
    hiding the syntactic and modeling complexity of each target from the programmer.
    \item To evaluate the resulting framework on CRDTs of increasing structural complexity, including recursive data
    structures, and characterize the cases in which automated and deductive verification agree or diverge.
\end{itemize}

\section{Document Structure}
\label{sec:structure}

\ntindex[Document Structure]{}

The remainder of this document is organized as follows:

\begin{itemize}
    \item Chapter~\ref{cha:related_work} presents the state of the art on consistency models, CRDT theory, and existing formal verification tools.
    \item Chapter~\ref{cha:background} details the theoretical foundations required for this work, including the formalization of CRDTs, how replicated data types are modeled and automatically verified in VeriFx, and the principles of deductive verification with Why3.
    \item Chapter~\ref{cha:language} presents the specification language introduced by the framework, its syntax, the
        \texttt{CvRDT} and \texttt{CmRDT} interfaces and the axioms attached to them, the constraints a module must
        satisfy to implement an interface, the directives that let a specification diverge between backends, and the
        type system underlying all of it.
    \item Chapter~\ref{cha:compiling} describes how a validated specification compiles to both VeriFx and Why3, the
        pattern each backend follows, how the compiler chooses between several generators depending on a CRDT's structure,
        and \texttt{RGA} as the case where every mechanism applies at once.
    \item Chapter~\ref{cha:case-studies} evaluates the framework across the ten CRDTs it specifies. It first compares the
        manual effort of writing each specification against the code generated for each backend, and then turns to the
        proofs themselves. Eight of the ten CRDTs are proved automatically in both VeriFx and Why3, and
        \texttt{PNCounter\_CRDT} needs a single manual step in Why3. For \texttt{RGA}, VeriFx times out, while Why3 proves
        that every pair of operations except two insertions commutes.
    \item Chapter~\ref{cha:conclusion} closes the dissertation, summarizing the problem addressed, the framework's
        design, and the evaluation results, and proposes extending the evaluation to further recursive or unbounded
        CRDTs as a direction for future work.
\end{itemize}
